\documentclass[10pt,a4paper]{article}
\usepackage[margin=0.78in]{geometry}
\usepackage{amsmath,amssymb,booktabs}
\usepackage[hidelinks]{hyperref}
\setlength{\parindent}{0pt}
\setlength{\parskip}{5pt}
\newcommand{\head}[1]{\par\medskip\textbf{#1}\par}
\newcommand{\Prule}{\mathcal P}
\begin{document}
\begin{center}
{\large\bfseries Native lane-flow equations: two admissibility checks}\\[3pt]
2026-10-20 --- analytical review, no traffic rollout
\end{center}

\textbf{Finding.} Literal source receiving adjustment can force negative
forward flow in a positive-density state. Separately, two pairwise lateral
desires can exceed their donor's common sending budget. These checks concern
the displayed equations, not an unobserved author implementation. They do
not establish a replacement rule, calibrated validity or a gradient result.

\head{1. Source rules and the exact interpretation checked}
Use the notation of Nagalur Subraveti et al.~\cite{ns}, with $L_{i,l\to l'}$
abbreviating their realized lateral flow $lq_{il\to l'}$. Density $k$ is in
veh/km and all $D,S,L,q$ are in veh/h. Their equations 6 and 13--17 give
\begin{align}
 \Prule_{i,l\to l'}&=\max\left\{0,
       \frac{I K_{il}-K_{il'}}{K_{il}+K_{il'}}\right\},\qquad
 D_{il}=\min\{u_{il}k_{il},C_{il}\}, \label{eq:desire}\\
 S_{il}&=\min\{C_{il},w_{il}(k^{\rm jam}_{il}-k_{il})\},\quad
 \phi_{i,l\to l'}=D_{il}\Prule_{i,l\to l'},\quad
 \theta_{il'}=\min\{1,(C_{il'}-S_{il'})/C_{il'}\}. \label{eq:receiving}
\end{align}
Here $K$ is a convex weighted average of current and two downstream densities
on that lane. For $\theta<1$, the printed realized flow is
$L=(1-\theta)\phi$. In both checks the targets have $S=C$, so $\theta=0$
and $L=\phi$; the disputed special branch at $\theta=1$ is never used.

For one lateral pair, source equations 20--22 print
\begin{align}
 D^{\rm mod}_{il}&=D_{il}+L_{i,l'\to l}-L_{i,l\to l'},\nonumber\\
 S^{\rm mod}_{il}&=S_{il}+L_{i,l'\to l}-L_{i,l\to l'},\quad
 q_{i-1,l}=\min\{D^{\rm mod}_{i-1,l},S^{\rm mod}_{il}\}. \label{eq:adjusted}
\end{align}
The second line reindexes the source receiving cell from $i+1$ to $i$.
It is not a sign correction. Thesis chapter 3, equations 3.20--3.22,
repeats this structure~\cite{thesis}.

For transparent arithmetic choose the illustrative triangular FD
\begin{equation}
 u=60\ {\rm km/h},\qquad w=20\ {\rm km/h},\qquad
 k^{\rm jam}=150\ {\rm veh/km},\quad k_c=37.5,\quad C=2250\ {\rm veh/h}.
 \label{eq:parameters}
\end{equation}
These are review examples, not fitted source-site parameters or a proposed
task calibration. The violations below are strict and tolerate small
parameter perturbations. One may use $\Delta x=0.1$ km and $\Delta t=6$ s,
satisfying the source CFL equality. No update trajectory is computed.

\head{2. One lateral departure makes printed receiving negative}
Consider one permitted direction from lane $l$ to a free target $l'$ away
from route/courtesy effects. Choose the source's allowed congested
keep-right setting zero, giving $I=1$. Other lateral flows are zero. Set
\[
\begin{array}{c|rrrr}
 \text{cell position}&i-1&i&i+1&i+2\\ \hline
 k_l&139&140&141&142\\
 k_{l'}&1&2&3&4
\end{array}
\]
All densities are positive and below jam density. At the cells needed here
they increase downstream, so the stated source discharge-mode condition
$k_i\geq k_{i+1}$ is inactive; the ordinary rules
\eqref{eq:desire}--\eqref{eq:receiving} apply. Since
$K_l\geq140$ and $K_{l'}\leq4$, the fraction is bounded by
\begin{equation}
 \Prule_{i,l\to l'}\geq\frac{140-4}{140+4}=\frac{17}{18},\qquad
 D_{il}=2250,\quad S_{il}=200,\quad S_{il'}=2250.
\end{equation}
Hence $L\geq2125$ and the printed receiving adjustment forces
\begin{equation}
 S^{\rm mod}_{il}=200-L\leq-1925,\qquad
 q_{i-1,l}\leq-1925\ {\rm veh/h}. \label{eq:negative}
\end{equation}
This needs no assumed upstream demand value: the minimum in
\eqref{eq:adjusted} cannot exceed its negative receiving argument. A
departure that should release space instead removes receiving allowance.
It is a forward-admissibility obstruction before VSL is added.

\head{3. Two pairwise desires can oversubscribe one donor}
Take a three-lane section away from structural incentives, with a middle
donor lane $l=2$. All three lanes are free. At $i,i+1,i+2$, use donor
densities $(30,31,32)$ and densities $(2,3,4)$ on each neighboring lane.
The current free-flow keep-right factor, where applied, is
$I=1-2/30=14/15$; an unrestricted direction has $I=1$, so $I\geq14/15$
for both outgoing directions. Reverse lateral desires are zero.
Consequently, independently of the convex lookahead weights,
\begin{equation}
 \Prule_{i,2\to j}\geq
 \frac{(14/15)30-4}{30+4}=\frac{12}{17},\qquad j\in\{1,3\}.
 \label{eq:two}
\end{equation}
Both targets have $S=C$; therefore with $D_{i2}=1800$,
\begin{equation}
 L_{i,2\to1}+L_{i,2\to3}\geq\frac{24}{17}D_{i2}
 =\frac{43200}{17}>1800\ {\rm veh/h}. \label{eq:overbudget}
\end{equation}
The source density update includes both lateral directions. If equation 20
is extended by summing both, as its conservation update requires for a
simultaneous two-direction interpretation, its residual forward demand is
\begin{equation}
 D^{\rm mod}_{i2}=D_{i2}-L_{i,2\to1}-L_{i,2\to3}
 \leq-\frac{12600}{17}<0. \label{eq:residual}
\end{equation}
Unlike the first check, this last step explicitly uses the natural summed
interpretation of the one-pair printed equation. If instead only one
direction is realized, a selection/priority rule must be supplied. The
displayed pairwise desires alone do not provide it. At the CFL equality,
the requested lateral masses alone total at least $72/17$ veh against a
current donor sending mass of 3 veh. This is oversubscription of its sending
budget, not a claim that a complete unobserved implementation loses mass.

Simply replacing $D$ by a positive VSL-capped budget while retaining the
two fractions gives the same ratio $24/17>1$. A cap does not allocate the
shared budget. Han et al.'s sending restriction~\cite{han} therefore cannot
alone resolve this issue.

\head{4. Later-source clarification and the decision boundary}
Thesis chapter 4, equation 4.11, uses unrestricted flow at $\theta=0$ and
$(1-\theta)\phi$ otherwise. On $0\leq\theta\leq1$ this is continuous and
gives zero at zero receiving. It differs from the original special branch
but neither hand check uses that branch. The later chapter also supplies
a three-to-two case. Neither observation resolves the original receiving
adjustment or supplies the joint allocator in the inspected summary.

An authoritative algorithm could clarify the transcription or allocation.
Absent that evidence, changing signs, clipping or selecting priorities
would create a changed flow-realization contract. Qualification stops at
that choice. No repaired law, replacement backbone, information advantage,
finite-clearance guarantee or derivative contract is selected here.

\begin{thebibliography}{3}\small
\bibitem{ns} H. H. S. Nagalur Subraveti, V. L. Knoop, and B. van Arem,
``First order multi-lane traffic flow model -- an incentive based macroscopic
model to represent lane change dynamics,'' \emph{Transportmetrica B}, vol. 7,
no. 1, pp. 1758--1779, DOI: \href{https://doi.org/10.1080/21680566.2019.1700846}{10.1080/21680566.2019.1700846}.
Source equations 6, 8--9, 13--17, 20--22; repository/copyright 2019,
article header 2020.
\bibitem{thesis} H. H. S. Nagalur Subraveti, \emph{Lane-specific Traffic Flow
Control}, Ph.D. dissertation, TU Delft, 2021,
\href{https://doi.org/10.4233/uuid:e405a549-b1fe-4fa1-8090-4429612359d9}{DOI: 10.4233/uuid:e405a549-b1fe-4fa1-8090-4429612359d9}.
Selected ch. 3, printed pp. 26--32; ch. 4, pp. 51--52.
\bibitem{han} Y. Han, A. Hegyi, Y. Yuan, S. Hoogendoorn, M. Papageorgiou,
and C. Roncoli, ``Resolving freeway jam waves by discrete first-order
model-based predictive control of variable speed limits,'' \emph{Transp.
Res. C}, vol. 77, pp. 405--420, 2017,
\href{https://doi.org/10.1016/j.trc.2017.02.009}{DOI: 10.1016/j.trc.2017.02.009}, sec. 2.3.
\end{thebibliography}
\end{document}
